\chapter{Tổng kết giai đoạn và kế hoạch phát triển}

\section{Tổng kết giai đoạn}

Giai đoạn 1 là giai đoạn đặt nền móng cho toàn bộ đồ án. Trong giai đoạn này, nhóm chưa triển khai hiện thực mã nguồn mà tập trung hoàn toàn vào hoạt động phân tích và thiết kế hệ thống. Việc đầu tư thời gian cho giai đoạn này giúp giảm thiểu rủi ro ở các giai đoạn sau, tránh tình trạng phải sửa đổi lớn do hiểu sai yêu cầu hoặc thiết kế không phù hợp.

Trước hết, nhóm đã tiến hành phân tích bài toán và bối cảnh nghiệp vụ, xác định rõ \textbf{mục tiêu tổng quát} của hệ thống, các \textbf{mục tiêu cụ thể} cũng như \textbf{phạm vi} của đồ án. Các đối tượng tham gia chính (người dùng cuối, quản trị viên, các dịch vụ tích hợp bên ngoài nếu có) được nhận diện rõ ràng, từ đó xác định các luồng tương tác quan trọng với hệ thống. Song song, nhóm cũng khảo sát các giải pháp tương tự đang tồn tại để tham khảo, rút kinh nghiệm về mặt chức năng và trải nghiệm người dùng.

Trên cơ sở đó, nhóm tiến hành thu thập và đặc tả yêu cầu. Các yêu cầu nghiệp vụ được chuyển hóa thành các yêu cầu chức năng và phi chức năng, bảo đảm mỗi yêu cầu đều có mô tả ngắn gọn, rõ ràng, có tiêu chí kiểm chứng. Những yêu cầu này được tổ chức thành các nhóm chức năng lớn (ví dụ: quản lý tài khoản, quản lý nội dung, tương tác giữa người dùng, quản trị hệ thống, \ldots) để thuận tiện cho việc thiết kế và triển khai sau này.

Trong giai đoạn phân tích, nhóm đã mô hình hóa nghiệp vụ và luồng xử lý thông qua các sơ đồ \textit{use-case}, \textit{activity} và \textit{sequence}. Các sơ đồ use-case giúp thể hiện góc nhìn của người dùng đối với hệ thống, nêu rõ hệ thống cần hỗ trợ những hành động nào. Các sơ đồ activity tập trung vào luồng xử lý của từng chức năng, nêu rõ điều kiện rẽ nhánh và các bước xử lý chính. Các sơ đồ sequence mô tả chi tiết sự tương tác theo thời gian giữa các đối tượng, là cơ sở quan trọng để thiết kế kiến trúc và giao diện lập trình ứng dụng (API).

Tiếp theo, nhóm đã tiến hành thiết kế kiến trúc tổng thể của hệ thống. Hệ thống được tổ chức theo kiến trúc phân lớp bao gồm lớp giao diện (frontend), lớp dịch vụ nghiệp vụ và API (backend) và lớp dữ liệu (database). Trong mỗi lớp, các module chính được xác định, ví dụ như module xác thực và phân quyền, module quản lý người dùng, module quản lý bài viết, module thông báo, module báo cáo, \ldots Kiến trúc cũng xác định rõ các giao tiếp giữa các module thông qua API, các chuẩn trao đổi dữ liệu (JSON, REST) và các cơ chế bảo mật cơ bản.

Thiết kế cơ sở dữ liệu là một phần quan trọng khác trong giai đoạn 1. Dựa trên các yêu cầu và mô hình nghiệp vụ, nhóm đã xây dựng mô hình thực thể -- kết hợp (ERD), sau đó chuyển đổi sang mô hình quan hệ chi tiết. Các bảng chính như người dùng, bài viết, tương tác, thông báo, đơn hàng (nếu có chức năng thương mại), \ldots được xác định cùng với các khóa chính, khóa ngoại và các ràng buộc toàn vẹn tham chiếu. Quá trình chuẩn hóa dữ liệu giúp loại bỏ dư thừa, đảm bảo dữ liệu nhất quán và dễ bảo trì về sau.

Song song với thiết kế kiến trúc và cơ sở dữ liệu, nhóm cũng xây dựng các bản thiết kế giao diện người dùng ở mức khung (wireframe). Các màn hình chính như trang đăng nhập, trang đăng ký, trang chủ, trang hồ sơ cá nhân, trang chi tiết bài viết, trang quản trị, \ldots được phác thảo, thể hiện bố cục, các thành phần giao diện cơ bản và luồng điều hướng giữa các màn hình. Các wireframe này là cơ sở quan trọng để thảo luận với giảng viên, người hướng dẫn và giữa các thành viên trong nhóm, giúp thống nhất sớm về trải nghiệm người dùng mong muốn.

Tổng thể, kết quả của giai đoạn 1 là bộ tài liệu phân tích và thiết kế tương đối đầy đủ, làm nền tảng cho việc hiện thực hệ thống ở giai đoạn 2. Nhờ có các tài liệu này, nhóm có được cái nhìn tổng thể, xác định rõ ràng các thành phần chính, luồng dữ liệu, giao tiếp giữa các module và các ràng buộc cần tuân thủ trong quá trình hiện thực.

\section{Đánh giá}

\subsection{Các kết quả đạt được}

Nhìn chung, giai đoạn 1 đã đạt được những kết quả tích cực, thể hiện qua các khía cạnh sau:

\begin{itemize}
    \item \textbf{Hoàn thiện tài liệu yêu cầu}: Hệ thống yêu cầu đã được mô tả rõ bằng use-case, kịch bản sử dụng và các ràng buộc phi chức năng. Mỗi yêu cầu đều gắn với đối tượng sử dụng và bối cảnh cụ thể, giúp dễ dàng kiểm chứng khi triển khai.
    \item \textbf{Thiết kế kiến trúc tương đối ổn định}: Kiến trúc phân lớp (frontend, backend, database) cùng các thành phần chính đã được xác định và liên kết hợp lý. Cách phân tách trách nhiệm giữa các lớp giúp hệ thống dễ mở rộng, dễ bảo trì và thuận tiện cho việc phân chia công việc giữa các thành viên.
    \item \textbf{Mô hình dữ liệu chặt chẽ}: Lược đồ cơ sở dữ liệu đã đảm bảo tính toàn vẹn, tránh dư thừa và hỗ trợ tốt cho các chức năng dự kiến. Các mối quan hệ một-nhiều, nhiều-nhiều giữa các thực thể được thể hiện rõ ràng, tạo thuận lợi cho việc truy vấn và tối ưu hóa hiệu năng.
    \item \textbf{Thiết kế giao diện nhất quán}: Các wireframe được xây dựng theo phong cách hiện đại, bố cục rõ ràng, tận dụng các thành phần giao diện quen thuộc với người dùng. Luồng điều hướng giữa các màn hình được thiết kế để giảm số bước thao tác, tăng tính trực quan.
    \item \textbf{Giảm rủi ro khi hiện thực}: Việc phân tích trước giúp dự đoán sớm các vấn đề về hiệu năng, bảo mật, mở rộng, từ đó có giải pháp thiết kế phù hợp (ví dụ: phân trang dữ liệu, cache, cơ chế phân quyền chi tiết). Nhờ vậy, khi bước vào giai đoạn hiện thực, nhóm ít phải thay đổi thiết kế ở mức lớn.
\end{itemize}

\subsection{Mức độ đáp ứng so với mục tiêu ban đầu}

So với mục tiêu đề ra khi bắt đầu đồ án, giai đoạn 1 cơ bản đã đáp ứng các yêu cầu chính:

\begin{itemize}
    \item Các \textbf{use-case cốt lõi} đều đã được xác định và mô hình hóa: đăng ký, đăng nhập, quản lý hồ sơ, đăng bài, tương tác với bài viết, quản trị nội dung, \ldots
    \item \textbf{Các sơ đồ phân tích} (use-case, activity, sequence, ERD, kiến trúc, \ldots) đều đã được xây dựng và liên kết với nhau, bảo đảm nhất quán về mặt khái niệm.
    \item \textbf{Phạm vi hệ thống} đã được làm rõ, giới hạn những chức năng nâng cao sẽ ưu tiên cho giai đoạn sau hoặc nằm ngoài phạm vi đồ án, giúp kế hoạch triển khai khả thi trong thời gian cho phép.
\end{itemize}

\subsection{Hạn chế và tồn tại}

Bên cạnh các kết quả tích cực, giai đoạn 1 vẫn còn một số hạn chế nhất định:

\begin{itemize}
    \item \textbf{Độ chi tiết một số mô hình còn hạn chế}: Một số use-case phức tạp, đặc biệt là các kịch bản lỗi, ngoại lệ hoặc xử lý song song, chưa được mô tả đủ chi tiết, cần bổ sung khi bước vào giai đoạn hiện thực. Điều này có thể dẫn đến việc phải điều chỉnh lại một số sơ đồ khi triển khai.
    \item \textbf{Chưa kiểm chứng bằng prototype}: Do chưa có mã nguồn hay prototype, tính khả thi về giao diện, trải nghiệm người dùng và hiệu năng mới chỉ dừng ở mức thiết kế lý thuyết. Một số quyết định thiết kế giao diện có thể cần điều chỉnh khi có phản hồi thực tế.
    \item \textbf{Ước lượng tài nguyên chưa thật chính xác}: Thời gian và nguồn lực cho một số chức năng nâng cao (ví dụ tích hợp AI, xử lý real-time, phân tích dữ liệu) còn mang tính tương đối. Khi bước vào giai đoạn hiện thực, nhóm có thể phải điều chỉnh mức độ ưu tiên hoặc phạm vi của các chức năng này.
    \item \textbf{Chưa có kiểm thử trên mô hình}: Mặc dù các sơ đồ đã được rà soát nội bộ, nhóm chưa có bước kiểm thử chính thức (ví dụ: walkthrough với người dùng giả định) để phát hiện các tình huống bất thường hoặc yêu cầu ẩn.
\end{itemize}

\subsection{Bài học kinh nghiệm}

Từ quá trình thực hiện giai đoạn 1, nhóm rút ra một số bài học kinh nghiệm quan trọng:

\begin{itemize}
    \item Cần phân rã yêu cầu và ưu tiên chức năng theo mức độ quan trọng và khả năng hiện thực, để lập kế hoạch tốt hơn và tránh dàn trải nguồn lực.
    \item Nên sớm xây dựng các prototype nhỏ cho giao diện hoặc luồng quan trọng (như đăng nhập, đăng bài, tìm kiếm) để kiểm chứng thiết kế, từ đó nhận phản hồi sớm và tinh chỉnh.
    \item Thường xuyên trao đổi giữa các thành viên frontend, backend, database để thống nhất giao tiếp (API, cấu trúc dữ liệu, định dạng trả về) ngay từ giai đoạn thiết kế, tránh mâu thuẫn khi triển khai.
    \item Cần quản lý phiên bản tài liệu một cách chặt chẽ, ghi lại lịch sử thay đổi để tiện truy vết khi cần tra cứu hoặc đối chiếu giữa các bản thiết kế khác nhau.
\end{itemize}

\subsection{Đề xuất cải tiến cho các giai đoạn sau}

Để giai đoạn 2 và các giai đoạn tiếp theo đạt hiệu quả cao hơn, nhóm đề xuất:

\begin{itemize}
    \item Bổ sung chi tiết cho các kịch bản use-case phức tạp, đặc biệt là các trường hợp lỗi, ngoại lệ, hành vi bất thường của người dùng.
    \item Xây dựng một số mockup hoặc prototype giao diện có thể thao tác được (clickable prototype) để kiểm chứng luồng tương tác và trải nghiệm người dùng.
    \item Xác định rõ các chỉ số kỹ thuật cần theo dõi (thời gian phản hồi, số lượng người dùng đồng thời dự kiến, dung lượng dữ liệu, \ldots) để có cơ sở đánh giá hiệu năng ở giai đoạn kiểm thử.
\end{itemize}

\section{Kế hoạch phát triển}

Giai đoạn 2 tập trung vào việc hiện thực hệ thống dựa trên thiết kế đã xây dựng, bao gồm: xây dựng cơ sở dữ liệu, hiện thực backend, frontend, tích hợp, kiểm thử và hoàn thiện tài liệu. Đây là giai đoạn tiêu tốn nhiều thời gian và nguồn lực nhất, đòi hỏi sự phối hợp chặt chẽ giữa các thành viên trong nhóm.

\subsection{Mục tiêu giai đoạn}

\begin{itemize}
    \item Hiện thực các chức năng cốt lõi theo yêu cầu đã phân tích, bảo đảm hệ thống có thể hoạt động end-to-end.
    \item Đảm bảo hệ thống hoạt động đúng, ổn định, có khả năng mở rộng và bảo trì, đồng thời tuân thủ các yêu cầu phi chức năng quan trọng (bảo mật, hiệu năng, khả dụng).
    \item Triển khai một phiên bản demo/phiên bản thử nghiệm để nghiệm thu với giảng viên, người hướng dẫn hoặc người dùng thử, thu thập phản hồi để hoàn thiện.
\end{itemize}

\subsection{Phạm vi và định hướng triển khai}

Trong giai đoạn 2, nhóm tập trung vào các nội dung sau:

\begin{itemize}
    \item Hoàn thiện cơ sở dữ liệu vật lý, viết script khởi tạo cấu trúc và dữ liệu mẫu phục vụ kiểm thử.
    \item Hiện thực các API backend theo tài liệu thiết kế, bao gồm các chức năng: xác thực, phân quyền, quản lý người dùng, quản lý nội dung, tương tác, thông báo, quản trị, \ldots
    \item Xây dựng giao diện frontend trên nền tảng công nghệ đã lựa chọn (ví dụ: React, Vite, Tailwind CSS), triển khai các màn hình và luồng tương tác theo wireframe.
    \item Tích hợp frontend và backend thông qua các API, xử lý các vấn đề về lỗi mạng, trạng thái tải dữ liệu, thông báo lỗi/thành công cho người dùng.
    \item Thực hiện kiểm thử chức năng, kiểm thử tích hợp, kiểm thử hệ thống và tinh chỉnh hiệu năng ở mức cơ bản.
\end{itemize}

\subsection{Kế hoạch chi tiết}

% \usepackage{longtable}

\begin{center}
\begin{longtable}{|c|p{5cm}|p{4cm}|p{3cm}|}
    \caption{Kế hoạch triển khai giai đoạn 2 trong 15 tuần} \label{tab:kehoach15tuan} \\
    \hline
    \textbf{Tuần} & \textbf{Nội dung công việc chính} & \textbf{Kết quả mong đợi} & \textbf{Ghi chú} \\
    \hline
    \endfirsthead

    \hline
    \textbf{Tuần} & \textbf{Nội dung công việc chính} & \textbf{Kết quả mong đợi} & \textbf{Ghi chú} \\
    \hline
    \endhead

    \hline
    \multicolumn{4}{r}{\textit{(trang sau)}} \\
    \hline
    \endfoot

    \hline
    \endlastfoot

    1 & Ôn lại tài liệu phân tích, rà soát thiết kế; thống nhất công nghệ, môi trường phát triển & Kế hoạch chi tiết, môi trường dev đã cấu hình & Họp nhóm, phân công cụ thể \\
    \hline
    2 & Thiết kế và tạo cơ sở dữ liệu vật lý, script khởi tạo dữ liệu mẫu, kiểm tra kết nối & CSDL hoạt động, có dữ liệu thử phục vụ phát triển và kiểm thử & \\
    \hline
    3 & Hiện thực các module backend nền tảng (xác thực, phân quyền, cấu trúc API, logging cơ bản) & API nền tảng chạy được, có kiểm thử đơn giản cho các chức năng chính & \\
    \hline
    4 & Hiện thực các chức năng nghiệp vụ cốt lõi (ví dụ: quản lý bài viết, người dùng) ở backend & Hoàn thiện API cho các use-case chính, có tài liệu mô tả endpoint & \\
    \hline
    5 & Mở rộng chức năng backend (thông báo, tin nhắn, báo cáo, \ldots), tối ưu truy vấn CSDL & Đa số API nghiệp vụ chính hoàn thành, hiệu năng ở mức chấp nhận được & \\
    \hline
    6 & Thiết kế và hiện thực khung giao diện frontend, cấu trúc routing, layout chính, component dùng chung & Ứng dụng frontend khung chạy được, kết nối thử với một số API mẫu & \\
    \hline
    7 & Hiện thực các màn hình chức năng cơ bản (đăng nhập, đăng ký, trang chủ, hồ sơ cá nhân) và luồng điều hướng & Frontend cho các chức năng cơ bản hoạt động với backend, xử lý được các lỗi phổ biến & \\
    \hline
    8 & Hiện thực các màn hình chức năng nâng cao (tìm kiếm, tương tác, quản lý nội dung, \ldots) và tối ưu UI/UX cơ bản & Hoàn thiện phần lớn giao diện và luồng chính, trải nghiệm người dùng nhất quán & \\
    \hline
    9 & Tích hợp toàn bộ hệ thống frontend--backend, xử lý lỗi, đồng bộ dữ liệu, tinh chỉnh thông báo cho người dùng & Hệ thống chạy end-to-end với dữ liệu thử, các lỗi tích hợp chính đã được xử lý & \\
    \hline
    10 & Kiểm thử chức năng, sửa lỗi, tối ưu hiệu năng ở mức cơ bản (cache, phân trang, giới hạn truy vấn) & Danh sách lỗi chính đã xử lý, hiệu năng đáp ứng yêu cầu ban đầu & \\
    \hline
    11 & Bổ sung chức năng phụ/trang quản trị (nếu có), tinh chỉnh UI/UX, chuẩn hóa thông điệp lỗi & Giao diện hoàn thiện hơn, trải nghiệm người dùng tốt hơn, hệ thống ổn định hơn & \\
    \hline
    12 & Kiểm thử hệ thống (system test), kiểm thử chấp nhận (UAT nội bộ), ghi nhận phản hồi & Báo cáo kiểm thử, danh sách lỗi còn lại và đề xuất cải tiến & \\
    \hline
    13 & Chuẩn bị dữ liệu demo, kịch bản trình diễn, tài liệu hướng dẫn sử dụng và cài đặt & Bộ dữ liệu demo, kịch bản demo và tài liệu hướng dẫn hoàn chỉnh & \\
    \hline
    14 & Hoàn thiện tài liệu báo cáo, cập nhật mô hình, sơ đồ theo phiên bản đã hiện thực, chuẩn bị slide & Bản nháp cuối của báo cáo và slide trình bày, thống nhất nội dung trình bày & \\
    \hline
    15 & Rà soát tổng thể, chỉnh sửa cuối cùng, chạy thử trình bày, nghiệm thu nội bộ trước khi bảo vệ & Phiên bản cuối cùng của hệ thống và báo cáo, sẵn sàng cho buổi bảo vệ & \\
    \hline
\end{longtable}
\end{center}

\subsection{Nguồn lực và phân công công việc}

Để kế hoạch trên khả thi, nhóm cần có sự phân công nhiệm vụ hợp lý giữa các thành viên. Thông thường, mỗi thành viên sẽ phụ trách chính một mảng (frontend, backend, cơ sở dữ liệu, kiểm thử/tài liệu), đồng thời hỗ trợ lẫn nhau khi cần. Việc phân công rõ ràng giúp tránh chồng chéo, đồng thời tăng tính trách nhiệm cá nhân.

Ngoài ra, nhóm cũng cần thống nhất các công cụ sử dụng (hệ quản trị cơ sở dữ liệu, IDE, hệ thống quản lý mã nguồn, công cụ quản lý công việc, \ldots) để đảm bảo quá trình phối hợp diễn ra trơn tru. Lịch họp định kỳ hàng tuần hoặc hai tuần một lần giúp cập nhật tiến độ, giải quyết sớm các vướng mắc phát sinh.

\subsection{Quản lý rủi ro và tiêu chí đánh giá}

Trong giai đoạn 2, một số rủi ro có thể xảy ra như: chức năng phức tạp hơn dự kiến dẫn đến trễ tiến độ, lỗi tích hợp khó xử lý, thay đổi yêu cầu muộn, thiếu hụt thời gian kiểm thử. Để hạn chế, nhóm cần thường xuyên đối chiếu tiến độ thực tế với kế hoạch 15 tuần, ưu tiên hoàn thành các chức năng cốt lõi trước, sau đó mới mở rộng thêm các chức năng phụ.

Tiêu chí đánh giá kết quả giai đoạn 2 không chỉ dừng ở việc “hoàn thành chức năng”, mà còn bao gồm: mức độ ổn định của hệ thống, khả năng đáp ứng yêu cầu phi chức năng, chất lượng tài liệu kèm theo, và mức độ hài lòng của người dùng thử hoặc giảng viên hướng dẫn.

\subsection{Sơ đồ Gantt}

% Cần thêm gói pgfgantt trong phần preamble của báo cáo:
% \usepackage{pgfgantt}

\begin{figure}[H]
    \centering
    \begin{ganttchart}[
        hgrid,
        vgrid,
        x unit=0.6cm,
        y unit chart=0.7cm
    ]{1}{15}
        \gantttitle{Kế hoạch triển khai giai đoạn 2 (15 tuần)}{15} \\
        \gantttitlelist{1,...,15}{1} \\

        % Phân tích và chuẩn bị
        \ganttgroup{Chuẩn bị}{1}{2} \\
        \ganttbar{Rà soát yêu cầu, thiết kế}{1}{1} \\
        \ganttbar{Cấu hình môi trường, CSDL}{2}{2} \\

        % Backend
        \ganttgroup{Hiện thực Backend}{3}{5} \\
        \ganttbar{API nền tảng (auth, role)}{3}{3} \\
        \ganttbar{Chức năng cốt lõi}{4}{4} \\
        \ganttbar{Chức năng mở rộng}{5}{5} \\

        % Frontend
        \ganttgroup{Hiện thực Frontend}{6}{8} \\
        \ganttbar{Khung giao diện, routing}{6}{6} \\
        \ganttbar{Màn hình cơ bản}{7}{7} \\
        \ganttbar{Màn hình nâng cao}{8}{8} \\

        % Tích hợp và kiểm thử
        \ganttgroup{Tích hợp và kiểm thử}{9}{12} \\
        \ganttbar{Tích hợp hệ thống}{9}{9} \\
        \ganttbar{Kiểm thử chức năng}{10}{10} \\
        \ganttbar{Hoàn thiện UI/UX}{11}{11} \\
        \ganttbar{Kiểm thử hệ thống, UAT}{12}{12} \\

        % Hoàn thiện và báo cáo
        \ganttgroup{Hoàn thiện và báo cáo}{13}{15} \\
        \ganttbar{Chuẩn bị dữ liệu demo, hướng dẫn}{13}{13} \\
        \ganttbar{Hoàn thiện báo cáo, slide}{14}{14} \\
        \ganttbar{Rà soát, nghiệm thu nội bộ}{15}{15}
    \end{ganttchart}
    \caption{Sơ đồ Gantt kế hoạch triển khai giai đoạn 2 trong 15 tuần}
\end{figure}


